\documentclass[10pt,compress]{beamer}
\mode<presentation>

\usepackage[english]{babel}
\usepackage[utf8]{inputenc}
\usepackage[T1]{fontenc}
\usepackage{lmodern}
\usepackage{listings}
\usepackage{multicol}

\usetheme{Madrid}
\setbeamertemplate{navigation symbols}{}
\setbeamertemplate{footline}[frame number]
\setbeamertemplate{caption}[numbered]

\title[Git]{Short Introduction to SCM, Git and GitHub}
\author{D Rodriguez, D F Barrero}
\institute{University of Alcala, UAH, Spain}
\date{}

\begin{document}

\begin{frame}
    \titlepage
\end{frame}

\begin{frame}{Objectives and References}
   \begin{block}{Objectives}
      \begin{enumerate}
         \item Understand why software projects need source control.
         \item Learn the basic Git workflow for day-to-day development.
         \item Identify the main collaboration features provided by GitHub.
            \item Understand the basics of CI/CD and security analysis in GitHub.
      \end{enumerate}
   \end{block}

   \begin{block}{References}
      \begin{enumerate}
          \item Git documentation: \url{https://git-scm.com/doc}
          \item GitHub Docs: \url{https://docs.github.com/}
          \item GitHub Skills: \url{https://skills.github.com/}
          \item Interactive tutorial: \url{https://learngitbranching.js.org/}
      \end{enumerate}
   \end{block}
\end{frame}

\begin{frame}[shrink]{Table of Contents}
 \begin{multicols}{2}
 \tableofcontents
 \end{multicols}
\end{frame}

\section{Software Configuration Management}

\subsection{Motivation}

\begin{frame}{Software Configuration Management}{Motivation}
\begin{itemize}
 \item Multiple developers edit the same codebase over time.
 \item Files evolve, bugs are fixed, and features are added incrementally.
 \item Teams need to know who changed what, when, and why.
 \item Reverting a broken change should be routine, not a crisis.
\end{itemize}

\bigskip

Without source control, collaboration becomes fragile and error-prone.
\end{frame}

\subsection{Version Control}

\begin{frame}{Software Configuration Management}{Version Control}
\begin{block}{Version control systems}
Version control systems keep track of changes to source code and project files.
They allow multiple people to edit a project in a predictable manner.
\end{block}

\begin{itemize}
 \item They store history.
 \item They help coordinate collaborative work.
 \item They make rollback and auditing possible.
\end{itemize}
\end{frame}

\subsection{Source Configuration Management}

\begin{frame}{Software Configuration Management}{Source Configuration Management}
Software configuration management is the task of tracking and controlling changes
in software artifacts, and it is part of the broader field of configuration management.

\bigskip

    Main open-source version control systems:
\begin{itemize}
 \item 1982: RCS
 \item 1990: CVS
 \item 2000: Subversion
 \item 2005: Git and Mercurial
\end{itemize}

There are many proprietary alternatives, but \texttt{Git} is now the dominant one.

\bigskip

    All software should live under version control. Otherwise, maintenance becomes risky.
\end{frame}

\section{Git}

\subsection{What is Git?}

\begin{frame}{Git}{What is Git?}
\begin{columns}
\column{.68\textwidth}
Git is an open source distributed version control system created by Linus Torvalds.

\bigskip

\url{https://git-scm.com/}

\bigskip

Compared with centralized systems, every clone contains the full repository history.

\column{.32\textwidth}
\begin{center}
 \includegraphics[width=.9\textwidth]{figs/torvalds-to-nvidia}
\end{center}
\end{columns}
\end{frame}

\subsection{Git Hosting Sites}

\begin{frame}{Git}{Git Hosting Sites}
It is usually easier to start with a hosted platform than to maintain your own server.

\begin{itemize}
 \item \alert{GitHub}: one of the most widely used Git hosting platforms; strong ecosystem and collaboration workflows.
 \item \alert{GitLab}: integrated DevSecOps platform with strong self-managed and cloud options.
    \item \alert{Bitbucket}: a Git hosting platform integrated with Atlassian tools such as Jira and Confluence.
\end{itemize}

Hosted platforms are commonly used as a central collaboration hub:
\begin{itemize}
 \item from which everyone pulls other people's changes
 \item to which everyone pushes their own changes
\end{itemize}

\tiny Note: plans and pricing details change frequently across providers.
\end{frame}

\begin{frame}{Git}{Git Hosting Platforms at a Glance}
\scriptsize
\begin{columns}[t]
\column{.33\textwidth}
\textbf{GitHub}
\begin{itemize}
\item Broad ecosystem and open source reach
\item GitHub Actions
\item Great default choice for many teams
\end{itemize}

\column{.33\textwidth}
\textbf{GitLab}
\begin{itemize}
\item Integrated DevSecOps platform
\item GitLab CI/CD
\item Strong self-managed options
\end{itemize}

\column{.33\textwidth}
\textbf{Bitbucket}
\begin{itemize}
\item Tight Jira and Confluence integration
\item Bitbucket Pipelines
\item Good fit for Atlassian-centered teams
\end{itemize}
\end{columns}
\end{frame}

\subsection{Git in IDEs}

\begin{frame}{Git}{Git in IDEs}
Most modern development environments integrate Git directly.

\begin{itemize}
 \item Commit, push, pull, merge, and resolve conflicts without leaving the editor.
 \item Inspect file history and diffs visually.
 \item Review branches, stashes, tags, and remotes from a GUI.
\end{itemize}

Examples include VS Code, PyCharm, IntelliJ IDEA, Eclipse, and many more.
\end{frame}

\subsection{Git vs. SVN}

\begin{frame}{Git}{Git vs. SVN (I)}
\begin{center}
\begin{columns}
\column{.50\linewidth}
\centering Centralized (SVN)\\\smallskip
\includegraphics[width=\linewidth]{figs/centralized.png}
\column{.50\linewidth}
\centering Distributed (Git)\\\smallskip
\includegraphics[width=\linewidth]{figs/distributed.png}
\end{columns}

\tiny \href{http://softwareengineering.stackexchange.com/questions/35074/im-a-subversion-geek-why-should-i-consider-or-not-consider-mercurial-or-git-or}{(Source)}
\end{center}
\end{frame}

\begin{frame}{Git}{Git vs. SVN (II)}
\begin{center}
\begin{columns}
\column{.50\linewidth}
\centering Fully distributed (Git)\\\smallskip
\includegraphics[width=\linewidth]{figs/fulldistributed.png}

\column{.50\linewidth}
\begin{block}{Git concepts to know}
\begin{itemize}
 \item \texttt{commit}
 \item \texttt{push}, \texttt{pull}, \texttt{fetch}
 \item \texttt{branch}, \texttt{merge}
 \item \texttt{origin}, \texttt{remote}
\end{itemize}
\end{block}
\end{columns}

\tiny \href{http://softwareengineering.stackexchange.com/questions/35074/im-a-subversion-geek-why-should-i-consider-or-not-consider-mercurial-or-git-or}{(Source)}
\end{center}
\end{frame}

\subsection{Repositories}

\begin{frame}{Git}{Local and Remote Repositories}
Key Git concepts to know:
\begin{itemize}
 \item \texttt{Working directory}: your project files as they currently exist on disk.
 \item \texttt{Repository}: the database that stores the history of changes.
 \item \texttt{Local repository}: the repository in your own machine.
 \item \texttt{Remote repository}: a copy of the repository hosted on a server.
 \item \texttt{Staging area}: the intermediary space used to prepare the next commit.
\end{itemize}
\end{frame}

\subsection{Workflow Overview}

\begin{frame}{Git}{Workflow Overview}
\centering
\includegraphics[width=.72\textwidth,height=.7\textheight,keepaspectratio]{figs/gitWorkflow.jpeg}
\end{frame}

\section{Using Git}

\subsection{Basic Workflow}

\begin{frame}{Using Git}{Basic Workflow}
Given initialized local and remote repositories, a common daily workflow is:

\begin{enumerate}
 \item Pull changes from the remote repository.
 \item Edit your files locally.
 \item Add the changes to the staging area.
 \item Commit the staged changes.
 \item Push your commits to the remote repository.
\end{enumerate}
\end{frame}

\subsection{Initializing a Repository}

\begin{frame}[fragile]{Using Git}{Initializing a Repository}
When using Git for the first time:

\begin{verbatim}
git config --global user.email user@example.com
git config --global user.name  "Jane Doe"
\end{verbatim}

Initialize a local repository:

\begin{verbatim}
mkdir /path/to/your/project
cd /path/to/your/project
git init
echo "# My project" > README.md
git add README.md
git commit -m "Initial commit"
git branch -M main
git remote add origin https://<where>/<project.git>
git push -u origin main
\end{verbatim}

\begin{itemize}
 \item \texttt{git init} creates the hidden \texttt{.git} directory.
 \item All Git metadata lives inside \texttt{.git}, which makes initialization safe.
 \item By default, a freshly initialized repository has no remotes.
\end{itemize}
\end{frame}

\begin{frame}{Using Git}{Cloning a Repository}
To work with an existing remote repository, the usual starting point is to clone it:

\bigskip

\texttt{git clone <repo>}

\bigskip

Example:

\texttt{git clone https://github.com/danrodgar/intro-git.git}

\bigskip

Once cloned, you can edit the repository locally as much as you want.
No changes reach the remote until you explicitly push them.
\end{frame}

\subsection{Basic Commands}

\begin{frame}[fragile]{Using Git}{Tracking Files}
To start tracking changes, we typically \texttt{add}, \texttt{commit}, and then \texttt{push}:

\begin{verbatim}
echo "A new file..." >> README.md
git add README.md
git commit -m "Initial commit"
git push -u origin main
\end{verbatim}

\begin{itemize}
 \item \texttt{git add} moves changes into the staging area.
 \item \texttt{git commit} records a snapshot of staged changes.
 \item \texttt{git push} sends your local commits to the remote repository.
\end{itemize}
\end{frame}

\begin{frame}{Using Git}{Pulling Changes}
Before starting new work, it is good practice to synchronize with the remote:

\begin{itemize}
 \item \texttt{git pull} downloads and integrates new remote changes.
 \item If you have unfinished local modifications, \texttt{git stash} can temporarily save them.
 \item \texttt{git status} shows the current state of your working tree.
 \item \texttt{git checkout -- <file>} or modern alternatives such as \texttt{git restore <file>} can discard local changes.
\end{itemize}
\end{frame}

\begin{frame}{Using Git}{Pushing Changes}
After committing locally, publish the result to the remote repository:

\bigskip

\texttt{git push origin main}

\bigskip

Pushing transfers your local commits. It does not automatically include uncommitted or unstaged changes.
\end{frame}

\subsection{Commits}

\begin{frame}[fragile]{Using Git}{Commits}
Each commit has:
\begin{itemize}
 \item an ID (its hash)
 \item an author
 \item a date
 \item a message describing the change
\end{itemize}

\begin{verbatim}
commit 44161dde6ea234f8cb997644f8e187123c3cc4af
Author: David <foo@foo.com>
Date:   Fri Mar 9 14:57:32 2018 +0100

    Issue with syntax highlighting solved
\end{verbatim}
\end{frame}

\subsection{Merge and Conflicts}

\begin{frame}[fragile]{Using Git}{Conflicts}
Merging is a common operation in Git.

\begin{itemize}
 \item Changes in different parts of a file are usually merged automatically.
 \item Conflicts happen when Git cannot decide how to merge two changes safely.
 \item Human intervention is then required.
\end{itemize}

\begin{columns}
\column{.50\textwidth}
\begin{block}{Conflict resolution}
\begin{enumerate}
 \item Identify the conflicted files.
 \item Open them and choose which changes to keep.
 \item Remove the conflict markers.
 \item Stage the resolved files and commit.
\end{enumerate}
\end{block}

\column{.40\textwidth}
\begin{exampleblock}{Conflict markers}
\begin{lstlisting}[basicstyle=\ttfamily\small]
def greet():
<<<<<<< HEAD
    return "Hello"
=======
    return "Hola"
>>>>>>> feature
\end{lstlisting}
\end{exampleblock}
\end{columns}
\end{frame}

\begin{frame}[fragile]{Using Git}{Handling Conflicts in Practice}
When a \texttt{git pull} reports conflicts, a common workflow is:

\begin{verbatim}
git status
# edit conflicted files and remove <<<<<<< ======= >>>>>>>
git add <resolved-file>
git commit
\end{verbatim}

If you want to cancel the merge attempt:
\begin{verbatim}
git merge --abort
\end{verbatim}

If you are resolving conflicts in a rebase:
\begin{verbatim}
git rebase main
# resolve files
git add <resolved-file>
git rebase --continue
\end{verbatim}
To cancel the rebase instead:
\begin{verbatim}
git rebase --abort
\end{verbatim}
\end{frame}

\begin{frame}{Using Git}{Merge and Conflicts: \texttt{diff}}
\texttt{diff -u <old file> <new file>}

\bigskip

This command shows the changes needed to turn the old file into the new file.

\begin{itemize}
 \item \texttt{---} and \texttt{+++} identify the old and new filenames.
 \item \texttt{@@} marks the location of the change.
 \item A leading space means the line is unchanged.
 \item \texttt{-} marks a deleted line.
 \item \texttt{+} marks an added line.
\end{itemize}
\end{frame}

\begin{frame}[fragile,shrink=5]{Using Git}{Merge and Conflicts: \texttt{diff} Example}
\begin{scriptsize}
\begin{columns}
\column{.45\textwidth}
\begin{block}{Old source file (\texttt{hello.c})}
\begin{verbatim}
#include <stdio.h>
int main() {
    printf("Hello World\n");
}
\end{verbatim}
\end{block}
\column{.45\textwidth}
\begin{block}{New source file (\texttt{hello\_new.c})}
\begin{verbatim}
#include <stdio.h>
int main(int argc, char *argv[]) {
  printf("Hello World\n");
  return 0;
}
\end{verbatim}
\end{block}
\end{columns}

\begin{alertblock}{Command used to generate the patch}
\begin{verbatim}
$ diff -u hello.c hello_new.c > hello.patch
\end{verbatim}
\end{alertblock}

\begin{exampleblock}{Generated patch output (\texttt{hello.patch})}
\begin{verbatim}
--- hello.c
+++ hello_new.c
@@ -1,4 +1,5 @@
 #include <stdio.h>
-int main() {
+int main(int argc, char *argv[]) {
   printf("Hello World\n");
+  return 0;
 }
\end{verbatim}
\end{exampleblock}
\end{scriptsize}
\end{frame}

\begin{frame}{Using Git}{Applying a Patch}
After generating a patch file:

\bigskip

\texttt{diff -u <old file> <new file> > file.patch}

\bigskip

you can apply it with:

\bigskip

\texttt{patch < file.patch}

\bigskip

The patch file contains enough information to identify the target file and the required edits.
\end{frame}

\begin{frame}{Using Git}{Original Patch}
\centering
\includegraphics[width=.4\textwidth,height=.75\textheight,keepaspectratio]{figs/patch.png}
\end{frame}

\subsection{Branches}

\begin{frame}{Using Git}{Branches}
\centering \includegraphics[width=.8\textwidth]{figs/branching}

\bigskip

\flushleft
Branches are used extensively to separate lines of development.

\begin{itemize}
 \item A repository can have multiple branches locally and remotely.
 \item The default branch is usually called \alert{main} or \alert{master}.
 \item \texttt{git branch <name>} creates a branch.
 \item \texttt{git checkout <branch>} or \texttt{git switch <branch>} changes the active branch.
 \item \texttt{HEAD} points to the current branch or, sometimes, to a specific commit.
\end{itemize}
\end{frame}

\subsection{Tags}

\begin{frame}{Using Git}{Tags}
A tag is a pointer to a specific point in repository history.

\begin{itemize}
 \item Tags usually have meaningful names such as \texttt{v1.0} or \texttt{release-2026-03}.
 \item They are commonly used to mark releases.
 \item Unlike branches, tags typically do not move after they are created.
\end{itemize}
\end{frame}

\subsection{Advanced Git}

\begin{frame}{Using Git}{Getting an Old Commit}
Sometimes you need to inspect an earlier state of the repository.

\begin{itemize}
 \item \texttt{git log}
 \item \texttt{git log --oneline}
\end{itemize}

These commands help locate old commits that you can inspect or recover files from.
\end{frame}

\subsection{Good Practices}

\begin{frame}{Using Git}{Good Practices}
Learn on the job: the best way to learn Git is by using it consistently.

\bigskip

Best practices:
\begin{itemize}
 \item Regularly push and pull.
 \item Test before pushing.
 \item Do not push half-baked changes.
 \item Do not pull in the middle of a risky unfinished task.
 \item Never commit temporary or generated files unless they belong in the repository.
 \item Keep commit messages short and informative.
 \item Keep the main branch stable and functional.
\end{itemize}

Remember: Git does not overwrite local changes without an explicit command.
\end{frame}

\section{GitHub}

\subsection{Features}

\begin{frame}{GitHub}{Features}
\begin{columns}
\column{.60\textwidth}
GitHub is a Git hosting platform with collaboration features on top of the repository itself.

\begin{itemize}
 \item Free public and private repositories (with plan limits)
 \item Repository browser
 \item Pull requests
 \item Issue tracking
 \item Code review workflows
 \item CI/CD with GitHub Actions
 \item Security analysis (CodeQL, Dependabot)
 \item Markdown rendering
 \item Project documentation and static hosting
\end{itemize}

\column{.40\textwidth}
\begin{center}
 \includegraphics[width=.75\textwidth]{figs/github-logo}
\end{center}
\end{columns}
\end{frame}

\begin{frame}{GitHub}{Key Concepts}
\begin{block}{Key GitHub concepts to know}
\begin{itemize}
 \item \texttt{Pull request}
 \item \texttt{Fork}
 \item \texttt{Issue}
 \item \texttt{README.md}
\end{itemize}
\end{block}

\begin{itemize}
 \item A fork creates your own copy of someone else's repository.
 \item A pull request proposes changes for review and discussion before merging.
\end{itemize}
\end{frame}

\begin{frame}{GitHub}{Pull Request Workflow}
Typical pull request (PR) flow:
\begin{enumerate}
 \item Create a branch from \texttt{main} (for example, \texttt{feature/login-form}).
 \item Commit focused changes and push the branch.
 \item Open a PR with clear title, context, and test notes.
 \item Request review and address comments with follow-up commits.
 \item Keep branch updated (merge or rebase from \texttt{main}).
 \item Merge when checks pass and review is approved.
\end{enumerate}

\begin{block}{Good PR habits}
Small PRs, descriptive commit messages, and clear review notes speed up integration.
\end{block}
\end{frame}

\subsection{Automation}

\begin{frame}{GitHub}{GitHub Actions}
GitHub Actions is GitHub's native workflow automation and CI/CD system.

\begin{itemize}
 \item Workflows are defined as YAML files in \texttt{.github/workflows}.
 \item Typical triggers: push, pull request, schedule, manual dispatch.
 \item Jobs can run tests, build artifacts, publish releases, or deploy.
 \item It integrates with pull requests to block merges when checks fail.
\end{itemize}

\begin{block}{Typical uses}
Test on each PR, lint code, build docs, run security checks, and deploy on tags.
\end{block}
\end{frame}

\subsection{Security}

\begin{frame}{GitHub}{CodeQL and Security Scanning}
CodeQL is GitHub's semantic code analysis engine for finding vulnerabilities and bugs.

\begin{itemize}
 \item Build a code database and query it with security-focused rules.
 \item Detects classes of issues that pattern-based linters can miss.
 \item Integrates with GitHub code scanning and pull request checks.
 \item Supports multiple languages (for example C/C++, Java, JavaScript, Python, and C\#).
\end{itemize}

\begin{block}{Reference}
\url{https://codeql.github.com/}
\end{block}
\end{frame}

\subsection{Repository Creation}

\begin{frame}{GitHub}{Repository Creation}
When creating a repository, you normally configure:
\begin{itemize}
 \item Name
 \item Description
 \item Visibility
 \item Initial \texttt{README.md}
 \item Optional \texttt{.gitignore} and license files
\end{itemize}

\begin{center}
 \includegraphics[width=.55\textwidth]{figs/newrepository.png}
\end{center}
\end{frame}

\subsection{README}

\begin{frame}{GitHub}{README}
Special file: \texttt{README.md}

\begin{itemize}
 \item Explains what the project is and how to use it.
 \item Is automatically rendered by GitHub.
 \item Usually includes installation, usage, and contribution instructions.
 \item The \texttt{md} extension stands for Markdown.
\end{itemize}
\end{frame}

\subsection{Markdown}

\begin{frame}[fragile]{GitHub}{Markdown}
Markdown is a lightweight markup language.

\begin{itemize}
 \item Simple to write
 \item Easy to read in plain text
 \item Powerful enough for documentation, READMEs, and web content
\end{itemize}

\bigskip

Typical elements include headings, emphasis, lists, links, images, and code blocks.
\end{frame}

\begin{frame}[fragile,shrink=8]{GitHub}{Markdown Example}
\begin{exampleblock}{Markdown example}
\begin{lstlisting}[basicstyle=\ttfamily\small]
# I am a header
## I am a subheader

Regular, *italic* and **bold**

- List item 1
- List item 2

[I am a link](http://foo.com)

![I am a pic](markdown.png)

~~~c
printf("Hello, world");
~~~
\end{lstlisting}
\end{exampleblock}
\end{frame}

\subsection{GitHub Pages}

\begin{frame}{GitHub}{GitHub Pages}
GitHub Pages integrates simple web publishing into the GitHub workflow.

\begin{itemize}
 \item Create project web sites and documentation portals.
 \item Publish from a dedicated branch or from a folder such as \texttt{docs}.
 \item It is commonly used together with Markdown-based content.
\end{itemize}

Typical project URL:

\url{https://<username>.github.io/<repository>}
\end{frame}

\end{document}
